\section{Results}
\label{sec:results}

The evidence is organized in three blocks. Section~\ref{sec:results:primary} tests the two conditions of operational usefulness (\emph{Specificity} and \emph{Matched-Control}) using a canonical intervention protocol. Section~\ref{sec:results:diagnostic} uses the same $41{,}784$ runs to analyze the two primary levers that move a pair from miss to hit (semantic-field composition and amplification scale) and two structural correlates of domain-level steerability (cross-prompt feature overlap and graph-scaffold compatibility). Section~\ref{sec:results:negative} records one within-domain null.

\begin{table*}[t]
\centering
\small
\caption{\textbf{Headline (C1): labeled supernodes vs.\ matched random controls.} For each domain we report (a) the per-pair best field-add variant with adaptive $M$-search Hit\% (a comprehensive operational measure), and (b) the labeled vs.\ matched-random vsMax discriminator at the full-labeled all-fields $M{=}20$ baseline (the cleanest specificity signal; full per-condition tables in \cref{appx:specificity}). Hit\% counts pairs whose first generated subword matches the target answer; vsMax is the best logit margin against the strongest competing dataset answer. C1 (\cref{def:opuseful}) is supported when both Hit-rate gap $\geq 5$~pp and vsMax gap $\geq 1.0$ logit. Sounds fails Condition~2 (vsMax gap $+0.14$) and is reported as a directional negative (\cref{sec:results:negative}).}
\label{tab:headline}
\begin{tabular}{lrrrrrrrl}
\toprule
& \multicolumn{3}{c}{Hit\% (best-FA, $M$-opt)} & & \multicolumn{3}{c}{vsMax (full lab.\ vs.\ rand., $M{=}20$)} & \\
\cmidrule(lr){2-4}\cmidrule(lr){6-8}
Domain & Lab. & Rand. & $\Delta$ & & Lab. & Rand. & $\Delta$ & Verdict \\
\midrule
USA       & 67.3 &  0.2 & $+67.1$ & & $+2.86$ & $-2.31$ & $+5.17$ & \textbf{C1 supported} \\
Books     & 63.8 &  0.0 & $+63.8$ & & $+5.98$ & $-0.15$ & $+6.13$ & \textbf{C1 supported} \\
Products  & 27.6 &  1.1 & $+26.5$ & & $+3.46$ & $+0.23$ & $+3.23$ & \textbf{C1 supported} \\
Paintings & 21.8 &  0.0 & $+21.8$ & & $+1.55$ & $-0.03$ & $+1.58$ & \textbf{C1 supported} \\
Sounds    & 20.0 &  3.3 & $+16.7$ & & $+3.28$ & $+3.14$ & $+0.14$ & \textit{Negative} (\cref{sec:results:negative}) \\
\bottomrule
\end{tabular}
\end{table*}


\subsection{Labeled supernodes have expected steering effect}
\label{sec:results:primary}

A supernode label is \emph{operationally useful} when both conditions hold: \textbf{(C1) Specificity}, the labeled intervention's absolute Hit-rate is at least $5$~pp and labeled vsMax is positive; and \textbf{(C2) Matched-Control}, the labeled vs.\ matched-random Hit-rate gap is at least $5$~pp and the vsMax gap is at least $1.0$ logit. We evaluate these conditions using the intervention harness of Section~\ref{sec:method:harness}. For each source--target pair $(e_A, e_B)$, the harness constructs an additive entity-swap on the concept-aligned supernodes of both entities and reports two operating points. The \emph{default} operating point fixes all three semantic fields (input, intermediate, answer) with ablation multiplier $M_{\mathrm{ablate}}{=}-2$ and amplification multiplier $M_{\mathrm{amplify}}{=}20$; the \emph{harness} operating point additionally selects the best field subset for the pair and accepts any amplification value identified by the adaptive $M$-search. The matched random control (\S\ref{sec:method:harness}) is evaluated at the same operating point as its labeled counterpart. Because labeled and random interventions share feature count and per-layer histogram by construction, any systematic gap between them reflects concept alignment.

At the default operating point (all-three-fields, $M{=}20$), the labeled--random vsMax gap is positive and exceeds $1$ logit in every domain (Table~\ref{tab:headline}), satisfying the vsMax component of both conditions. Labeled Hit-rate also exceeds matched random everywhere at this default, but in Books and Paintings the absolute labeled Hit-rate sits below the $5$~pp Specificity threshold. At the harness operating point (best field subset plus adaptive $M$), labeled Hit-rate rises into the tens of percent while matched random stays below $2$\% in every domain, clearing the absolute threshold in all four. Per-condition tables and per-pair confidence intervals are in \cref{appx:specificity}.

Suppression of the source answer happens easily: matched random ablation suppresses the source as effectively as the labeled intervention in three of four domains (\cref{appx:specificity}). What separates the two conditions is \emph{targeted redirection}---only the labeled intervention concentrates in a clean-redirection regime where the target overtakes the source at position~$0$, while matched random concentrates in a generic-disruption regime that promotes nothing in particular (\cref{appx:regime}). The discriminating signal is therefore redirection toward the target, not removal of the source.

All four domains satisfy both Specificity and Matched-Control (Table~\ref{tab:headline}): the vsMax gap clears $1.0$ logit at the default operating point, and both the absolute Hit-rate and the labeled--random Hit-rate gap clear $5$~pp at the harness operating point.

\subsection{Levers and correlates of steerability}
\label{sec:results:diagnostic}

The analyses below use the same $41{,}784$ runs to characterize variation in intervention quality. Each is a follow-up to the primary test of Section~\ref{sec:results:primary}; the four-domain verdict above does not depend on any of these diagnostics individually.

\subsubsection{Intermediate and answer fields carry the steering signal}
\label{sec:results:diagnostic:fields}
Restricting the intervention to a two-field subset substantially improves redirection over the default all-three-field configuration. In USA the best subset (\emph{state + capital}) raises Hit\% from $24.7$ to $38.8$; in Books (\emph{book + author}), from $3.8$ to $37.1$; in Products (\emph{company + founder}), from $15.2$ to $24.2$. Paintings has a weaker baseline and a correspondingly weaker gain (best subset \emph{first\_name}, $3.3$ to $6.7$). Aggregated across the four domains, the intermediate field alone drives Hit\% $14.6\%$, the answer field $8.8\%$, and the input field only $6.1\%$ (\cref{appx:fieldadd}). Input-field features appear to encode the model's representation of the surface form, while the circuits that select the emitted answer concentrate in the intermediate and answer fields; including input-field features dilutes the redirection signal with activity that is not aligned with the target answer.

\subsubsection{$M{=}20$ overshoots in the multi-feature regime}
\label{sec:results:diagnostic:msearch}
The $M{=}20$ amplification default inherited from single-feature steering overshoots in the multi-feature regime. Adaptive $M$-search rescues $447$ pairs across all conditions, and the distribution of winning values is bimodal, with peaks near $M{\sim}2.4$ ($47$\% of labeled hits) and $M{\sim}6.9$ ($40$\%), both below the default (\cref{appx:msearch}). Two controls rule out the alternative that the default is simply too small. Testing $M \in \{50, 100, 200\}$ on $80$ near-miss pairs across three domains produced $0$ additional hits; KL divergence between the steered and unsteered distributions saturates above $M{\sim}50$ while the generated text continues to degrade. Testing the top-$k$ features by graph influence ($k \in \{1, 3, 5, 10\}$) at $M \in [20, 200]$ on Products produced $0$ hits out of $96$, and only the full feature set yields hits. The steering signal is distributed across the feature set, and multi-feature amplification at high $M$ linearly amplifies every direction in the labeled set simultaneously, producing distributed noise that dilutes the target signal. The single-feature SAE setting of \citet{templeton2024scaling}, in which clamping at roughly $\pm 10\times$ maximum observed activation saturates one direction, behaves differently for a structural reason: a single direction cannot linearly superpose with itself.

\subsubsection{Shared early-layer backbone, entity-specific late layers}
\label{sec:results:diagnostic:layers}
A scalable cross-prompt analysis on $1{,}592$ intra-domain entity pairs spanning $97$ entities characterizes how each entity's feature set distributes across the $26$ layers of Gemma-2-2B-it (\cref{appx:crossprompt}). Activation stability---the probability that a feature active on the seed prompt is also active on a concept-matched probe for the same entity---exceeds $90$\% in all domains, so the CPAS picture of each entity is reproducible at scale. Jaccard overlap between entity feature sets within a domain ranges from $0.29$ to $0.47$, all significantly above chance under permutation ($p < 0.001$ at several pool sizes), and the ratio of early-layer to late-layer Jaccard is between $1.4$ and $3.0\times$ in three of four domains. The pattern is consistent with a shared backbone of early-layer features reused across entities and entity-specific specialization concentrated in later layers.

\subsubsection{Scaffold compatibility predicts cross-domain steerability}
\label{sec:results:diagnostic:scaffold}
Scaffold compatibility tracks cross-domain steerability. We define per-domain \emph{scaffold compatibility} as the fraction of total graph influence carried by features that receive identical supernode assignment across two entity graphs. With one scaffold number per domain, the metric orders the four domains by Hit\% at the default operating point in the same sequence (Spearman $\rho{=}1.0$, $N{=}4$; \cref{appx:scaffold}): USA (scaffold $0.53$, Hit\% $24.7$), Products ($0.42 \rightarrow 15.2$), Paintings ($0.36 \rightarrow 4.4$), and Books ($0.25 \rightarrow 3.8$). Restricting to later layers ($\geq L_{16}$) makes the gradient steeper. The cross-domain ordering does not carry down to a pair-level predictor within a single domain: in USA, splitting the $2{,}450$ pairs into quartiles of pair-level scaffold compatibility gives nearly flat Hit-rates ($27.8$, $24.7$, $24.3$, $24.3$\%) with absolute Pearson correlations below $0.10$ for every variant tested, plausibly because the within-USA scaffold range is narrow ($[0.38, 0.73]$). In smaller domains the within-domain signal is weak but non-zero (e.g., residualized $r{=}+0.29$ in Products with $N{=}132$). We therefore limit the scaffold claim to the between-domain level and report the within-USA null in Section~\ref{sec:results:negative}.

Collectively, these diagnostics account for the gap between the default and harness operating points. Semantic-field composition and amplification scale are the two levers that move a pair from miss to hit under the same labeled feature set, and the cross-prompt overlap and scaffold analyses locate a structural correlate of domain-level steerability.

\subsection{Scaffold compatibility does not predict within-USA Hit-rate}
\label{sec:results:negative}

We record one within-domain null alongside the primary evaluation: the scaffold metric does not predict pair-level Hit-rate within USA, our most statistically powered domain. Splitting the $2{,}450$ USA pairs into four quartiles of pair-level scaffold compatibility gives Hit-rates of $27.8 / 24.7 / 24.3 / 24.3\%$, essentially flat, with $|r|<0.10$ for all variants tested. The likely reason is that within-domain scaffold values in USA lie in the narrow range $[0.38, 0.73]$, too narrow to separate pairs. The scaffold claim is therefore between-domain only, and we do not use scaffold compatibility as a within-domain predictor.
